\documentclass[12pt]{jlreq}
\usepackage{amsmath, amssymb}

\newcommand{\C}{\mathbb{C}}
\newcommand{\R}{\mathbb{R}}
\newcommand{\Z}{\mathbb{Z}}
\newcommand{\N}{\mathbb{N}}
\newcommand{\F}{\mathbb{F}}
\newcommand{\Prob}{\mathbb{P}}
\newcommand{\Exp}{\mathbb{E}}
\newcommand{\dd}{\,d}
\newcommand{\Ker}{\operatorname{Ker}}
\newcommand{\Impart}{\operatorname{Im}}
\newcommand{\Hom}{\operatorname{Hom}}

\begin{document}

\(q\) を \(0<|q|<1\) をみたす複素数とする。\([0]:=1,\ [0]!:=1\)，正の整数 \(n\) に対して
\[
  [n]:=\frac{1-q^n}{1-q},\qquad
  [n]!:=[n][n-1]\cdots[1],\qquad
  \begin{bmatrix}n\\ k\end{bmatrix}:=\frac{[n]!}{[n-k]![k]!}
\]
と定義する。また，複素平面内の領域 \(X:=\{x\in\C\mid |x|<d\}\ (d>0)\) で定義された関数 \(f(x)\) に対して，
\[
  Bf(x):=
  \begin{cases}
    \dfrac{f(x)-f(qx)}{(1-q)x} & (x\ne0),\\
    0 & (x=0),
  \end{cases}
\]
\[
  B^0f(x):=f(x),\qquad
  B^{k+1}f(x):=B(B^kf(x))\qquad (k=0,1,2,\ldots)
\]
とする。

次に，\(a_k,\ b_k\ (k=0,1,2,\ldots)\) を複素定数とし，べき級数
\[
  \psi(x):=\sum_{k=0}^{\infty}b_kx^k,\qquad
  \Omega(x):=\sum_{k=0}^{\infty}a_kb_kx^k
\]
はともに，\(X\) で絶対収束するものとする。そして，\(\{a_k\}_{k=0}^{\infty}\) に作用する演算 \(E\) および \(D_r\ (r=0,1,2,\ldots)\) を
\[
  Ea_k=a_{k+1}\qquad (k=0,1,2,\ldots),
\]
\[
  D_0:=1,\qquad
  D_r:=(E-1)(E-q)\cdots(E-q^{r-1})\qquad (r=1,2,3,\ldots)
\]
によって定義する。たとえば
\[
  D_0a_0=a_0,\qquad
  D_1a_0=a_1-a_0,\qquad
  D_2a_0=a_2-(1+q)a_1+qa_0
\]
である。以下の問に答えよ。

\begin{enumerate}
  \item 任意の正整数 \(k\) に対して，
  \[
    E^k=\sum_{j=0}^{k}\begin{bmatrix}k\\ j\end{bmatrix}D_j
  \]
  が成り立つことを示せ。
  \item \(x\in X\) に対して，次の等式が成り立つことを示せ。ただし，右辺の級数の収束は仮定してよい。
  \[
    \Omega(x)=\sum_{k=0}^{\infty}\frac{(D_ka_0)x^k}{[k]!}B^k\psi(x).
  \]
  \item \(\displaystyle\psi(x)=\frac{1}{(1-x)(1-qx)}\) を考えることによって，\(|x|<1\) であるとき，
  \[
    \sum_{k=0}^{\infty}[k+1]^2x^k=
    \frac{1+qx}{(1-x)(1-qx)(1-q^2x)}
  \]
  であることを示せ。
\end{enumerate}

\end{document}